\section{Observabilidad y prácticas de depuración}
\subsection{Registro y pipeline de monitoreo}
Para que cada ejecución del Agent se pueda rastrear, OpenHands abstrae la terminal, el navegador, las pruebas y los registros en un flujo de eventos unificado. Cada evento incluye una marca de tiempo, la herramienta invocada, una instantánea del contexto y el resultado esperado, lo que facilita a los ingenieros localizar el punto de la falla.

\begin{itemize}
    \item Registro de terminal: registra cada comando de shell y su código de retorno, y se sube junto con los resultados de las pruebas.
    \item Interacción de navegador y red: captura clics, envíos de formularios y solicitudes/respuestas, lo que facilita reproducir sesiones HTTP y flujos de API.
    \item Diff de archivos/código: guarda automáticamente los fragmentos de diff modificados por el Agent, las rutas de los archivos relacionados y el contexto.
\end{itemize}

\begin{table}[H]
\centering
\begin{tabular}{p{0.25\textwidth} p{0.35\textwidth} p{0.33\textwidth}}
\toprule
Nivel & Uso & Ejemplo \\
\midrule
Nivel de comando & Captura cada operación de shell y su resultado & \texttt{pip install}, \texttt{pytest tests/foo.py} \\
\midrule
Nivel de contexto & Registra la versión del contexto usada y los datos recuperados & commit de Git + Issue summary correspondientes \\
\midrule
Nivel de resultado & Archiva si las pruebas pasaron o fallaron, análisis de diff & Unit test failed with stack trace \\
\bottomrule
\end{tabular}
\caption{Niveles de observación y usos del registro en OpenHands.}
\end{table}

\subsection{Modo de depuración y reversión rápida}
Cuando el Agent planea ejecutar un comando de alto riesgo (como una publicación o una migración de base de datos), el sistema entra automáticamente en modo de depuración: pausa el proceso, captura la instantánea actual y le pide a un humano que confirme el comando. Si la ejecución del comando falla, la plataforma puede revertirla automáticamente con base en los registros y el diff, evitando contaminar la rama principal.

\begin{knowledgebox}{Recomendación de depuración}
Hacer que el Agent registre en el log la «intención» y el «resultado esperado» de cada paso puede acelerar la revisión manual y formar muestras de entrenamiento.
\end{knowledgebox}

\begin{warningbox}{El uso indebido de los registros puede filtrar información sensible}
Los registros de depuración pueden contener claves de API, rutas y trazas de pila, por lo que es necesario establecer anonimización y control de acceso para evitar filtraciones de privacidad durante las auditorías.
\end{warningbox}

\subsection{Resumen de la sección}
\begin{itemize}
    \item Con un flujo de eventos unificado y registros con plantilla se puede acotar rápidamente el alcance de las fallas.
    \item El modo de depuración ofrece una doble garantía: confirmación manual y reversión automática.
\end{itemize}
